\chapter{Statistical Significance of Model Comparisons}
\label{app:significance}

This appendix reports the statistical tests behind the model comparisons in
Section~\ref{ssec:multi-frame Fusion}. Three fusers are compared on the static
test set: the recurrent baseline \textbf{MISRGRU}~\cite{resurf2024}, the
transformer fuser \textbf{TR-MISR} (8-frame), and the
\textbf{DPA} model (streaming MISRGRU with the Deformable Phase-space Alignment
module.

\section*{Procedure}
\label{appsec:sig_method}

All comparisons are paired. Every model is evaluated on the identical
480-sample test set (12 blinking-speed $\times$ photon-level conditions,
40 samples each), and scores are matched by condition and sample identity. For a
pair of models $A$ and $B$ we form the per-sample differences of the
resolution-scaled Pearson correlation score, $d_i = \mathrm{RSP}^A_i - \mathrm{RSP}^B_i$, and
test the null hypothesis that their median is zero.

Because RSP is bounded in $[-1,1]$ and not normally distributed, the primary test
is the two-sided \textbf{Wilcoxon signed-rank test}. When several pairs are
compared within one family of conditions, the $p$-values are adjusted for
multiple comparisons with the \textbf{Holm--Bonferroni} method, and
significance is declared at an adjusted level of $\alpha = 0.05$.

We also show the \textbf{mean RSP difference}
$\overline{d}$, and
the \textbf{matched-pairs rank-biserial correlation} $r \in [-1, +1]$ measures how
consistently one model wins ($r = +1$ means $A$ beats $B$ on every sample,
$r = 0$ means wins and losses balance). Because the sample sizes are large
($n = 40$ per condition, $n = 480$ pooled), the $p$-values are extremely small and
several saturate the exact-test floor; interpretation therefore relies on the
effect sizes rather than on the $p$-values themselves.

\section*{Three-model comparison at a 20-frame budget}
\label{appsec:sig_20f}

At a matched budget of 20 input frames we compare all three models pooled over
the full test set ($n = 480$).

\begin{table}[htbp]
  \centering
  \caption{Paired comparison of the three fusers at a 20-frame budget, pooled over
           all 480 test samples. $\overline{d}$ is the mean RSP difference
           (A\,$-$\,B); $p$ is Holm-adjusted Wilcoxon. Higher RSP is better.}
  \label{tab:sig_20f}
  \begin{tabular}{lccccc}
    \toprule
    \textbf{Comparison (A vs.\ B)} & \textbf{mean A} & \textbf{mean B} &
    $\overline{d}$ & \textbf{better} & \textbf{rank-biserial} \\
    \midrule
    MISRGRU vs.\ TR-MISR & 0.8895 & 0.9025 & $-0.0129$ & TR-MISR & $-0.81$ \\
    MISRGRU vs.\ DPA     & 0.8895 & 0.9024 & $-0.0128$ & DPA     & $-0.86$ \\
    TR-MISR vs.\ DPA     & 0.9025 & 0.9024 & $+0.0001$ & - & $-0.06$ \\
    \bottomrule
  \end{tabular}
\end{table}

Both TR-MISR and DPA significantly outperform the MISRGRU baseline
($p \approx 3.1\times10^{-52}$ and $8.2\times10^{-60}$ respectively), with large
and highly consistent effects (rank-biserial $\approx -0.8$; MISRGRU loses on the
large majority of samples). The final streaming MISRGRU
(Section~\ref{sec:streaming}) was evaluated with the same paired procedure and
likewise beats the batch MISRGRU baseline at a matched 20-frame budget
($\overline{d} = 0.010$ RSP in its favour, rank-biserial $-0.99$,
$p \approx 6\times10^{-78}$). It wins on almost every paired sample. In contrast, TR-MISR and DPA are
statistically indistinguishable in aggregate: the mean difference is
$+0.0001$ RSP and the Wilcoxon test is not significant ($p = 0.24$).

This aggregate tie is misleading. When we compare the models per imaging setting it
reveals a clean difference: DPA is significantly better in the fast-blinking /
high-SNR conditions ($\overline{d}$ from $-0.015$ to $-0.024$ in favour of DPA),
whereas TR-MISR is significantly better in the slow-blinking / low-SNR
conditions, its advantage growing as the signal weakens (up to $+0.040$ RSP at
the hardest 12\_slow\_p40 condition). 

\section*{DPA vs.\ TR-MISR at a 30-frame budget}
\label{appsec:sig_30f}

Increasing the number of input frames to 30 frames for both models confirms the same
pattern. Pooled over all conditions ($n = 480$), DPA holds a small but significant
edge ($\overline{d} = -0.0034$ RSP, $p = 2.9\times10^{-6}$, rank-biserial
$-0.25$), a tie that hides a tradeoff.

\begin{table}[htbp]
  \centering
  \caption{DPA vs.\ TR-MISR at a 30-frame budget, grouped by blinking speed.
           $\overline{d}$ is the mean RSP difference (TR-MISR\,$-$\,DPA);
           $\star$ marks Holm-adjusted Wilcoxon $p < 0.05$.}
  \label{tab:sig_30f}
  \begin{tabular}{lcccc}
    \toprule
    \textbf{Condition group} & $\overline{d}$ \textbf{(TR-MISR\,$-$\,DPA)} &
    \textbf{better} & \textbf{rank-biserial} & \textbf{sig} \\
    \midrule
    fast (p100--p40)   & $-0.019$ to $-0.023$ & DPA     & $-1.00$            & $\star$ \\
    medium (p100--p60) & $-0.006$ to $-0.012$ & DPA     & $-0.99$ / $-1.00$  & $\star$ \\
    medium p40         & $+0.001$             & --- (n.s.) & $+0.30$          &         \\
    slow (p100--p40)   & $+0.008$ to $+0.032$ & TR-MISR & $+1.00$            & $\star$ \\
    \bottomrule
  \end{tabular}
\end{table}

The sign of the difference flips monotonically with blinking speed: under fast
blinking DPA wins every sample (rank-biserial $=-1.0$), the medium conditions form
the crossover region, and under slow blinking TR-MISR wins every sample, with its
margin increasing as the photon count drops (up to $+0.032$ RSP at
12\_slow\_p40). The photon level does not change which model wins but modulates the
size of the gap.

In summary, at equal frame budgets both fusers clearly beat the recurrent MISRGRU
baseline, but neither dominates the other: DPA is the stronger model under
fast-blinking / high-SNR conditions, TR-MISR under slow-blinking / low-SNR
conditions, the two crossing over in the medium regime. Because these opposing
effects cancel when conditions are pooled, the comparison is only meaningful when
stratified.
